\documentclass[12pt,oneside]{report}
\usepackage[T1]{fontenc}
\usepackage[utf8]{inputenc}
\usepackage[ngerman]{babel}
\usepackage[sc]{mathpazo}		% Palatino für Haupttext und Mathematik
\usepackage{helvet}
\usepackage{microtype}
\usepackage{parskip}			% Absätze mit Abstand statt Einrückung

\usepackage[a4paper,
	    left=2.5cm,
	    right=2.5cm,
	    top=1.5cm,
	    bottom=1.5cm,
	    marginparsep=5mm,
	    marginparwidth=10mm,
	    headheight=7mm,
	    headsep=1.2cm,
	    footskip=1.5cm,
	    includeheadfoot]{geometry}

\usepackage{amsmath,amssymb}
\usepackage{graphicx}
\usepackage{booktabs}
\usepackage{tabularx}
\usepackage{xcolor}
\usepackage{tcolorbox}
\tcbuselibrary{skins,breakable}

% Farbdefinitionen für Links und Warnhinweise
\definecolor{linknavy}{RGB}{0, 60, 130}

% Box-Farben für Warnhinweise und Infoboxen
\definecolor{warnred}{RGB}{180, 40, 40}
\definecolor{warnbg}{RGB}{253, 242, 242}
\definecolor{infoblue}{RGB}{0, 90, 160}
\definecolor{infobg}{RGB}{240, 247, 255}

%%%%%%%%%%%%%%%%%%%%%%%%%%%%%%%%%%% Code Listings %%%%%%%%%%%%%%%%%%%%%%%%%%%%%%%%%%%
\usepackage{listings}

% VS Code Light Farben (wie in master.tex)
\definecolor{vscBackground}{HTML}{FFFFFF}
\definecolor{vscText}{HTML}{000000}
\definecolor{vscKeyword}{HTML}{0000FF}
\definecolor{vscString}{HTML}{A31515}
\definecolor{vscComment}{HTML}{008000}
\definecolor{vscNumber}{HTML}{098658}
\definecolor{vscLineNumber}{HTML}{237893}

\lstdefinestyle{vscodeLightMatlab}{
    language=Matlab,
    backgroundcolor=\color{vscBackground},
    basicstyle=\ttfamily\footnotesize\color{vscText},
    keywordstyle=\color{vscKeyword},
    stringstyle=\color{vscString},
    commentstyle=\color{vscComment},
    numberstyle=\tiny\color{vscLineNumber},
    numbers=left,
    numbersep=8pt,
    breaklines=true,
    breakatwhitespace=false,
    keepspaces=true,
    showspaces=false,
    showstringspaces=false,
    showtabs=false,
    tabsize=4,
    captionpos=b,
    literate=
        {ä}{{\"a}}1
        {ö}{{\"o}}1
        {ü}{{\"u}}1
        {Ä}{{\"A}}1
        {Ö}{{\"O}}1
        {Ü}{{\"U}}1
        {ß}{{\ss}}1
}

\lstset{style=vscodeLightMatlab}
%%%%%%%%%%%%%%%%%%%%%%%%%%%%%%%%%%%%%%%%%%%%%%%%%%%%%%%%%%%%%%%%%%%%%%%%%%%%%%%%%%%%%

\newcommand{\DocVersion}{1.1}

\usepackage{fancyhdr}
\pagestyle{fancy}
\fancyhf{}
\setlength{\headwidth}{\textwidth}
\fancyfoot[R]{\thepage}
\fancyfoot[L]{\leftmark}
\fancyhead[R]{SES Panel Suite -- User Guide v\DocVersion}
\fancyhead[L]{Formula Student SES}
\renewcommand{\chaptermark}[1]{\markboth{Kapitel \thechapter. #1}{}}
\renewcommand{\headrulewidth}{0.5pt}
\renewcommand{\footrulewidth}{0.5pt}

\fancypagestyle{plain}{
  \fancyhf{}
  \fancyfoot[R]{\thepage}
  \fancyfoot[L]{\leftmark}
  \fancyhead[R]{SES Panel Suite -- User Guide v\DocVersion}
  \fancyhead[L]{Formula Student SES}
  \renewcommand{\headrulewidth}{0.5pt}
  \renewcommand{\footrulewidth}{0.5pt}
}

\usepackage[bookmarks, raiselinks, pageanchor,
            hyperindex, colorlinks,
            citecolor=black, linkcolor=black,
            urlcolor=blue!70!black, filecolor=black,
            menucolor=black]{hyperref}

\hypersetup{
    pdftitle={SES Panel Calculation and Optimization Suite -- User Guide v\DocVersion},
    pdfauthor={Formula Student Team -- Chassis and SES Department},
    pdfsubject={User Guide: SES Sandwich-Panel-Berechnung},
    pdfkeywords={Formula Student, SES, CFK, Sandwich, Wabe, 3-Punkt-Biegung, Optimierung}
}

% Makros für Verlinkungen auf das Abkürzungsverzeichnis (\label{abk:...})
\newcommand{\abk}[1]{\hyperref[abk:#1]{#1}}
\newcommand{\abklink}[2]{\hyperref[abk:#1]{#2}}

\begin{document}
  \hypersetup{pageanchor=false}

  % ---------------------------------------------------------------------------------
  % Titelseite
  % ---------------------------------------------------------------------------------
  \begin{titlepage}
    \thispagestyle{empty}
    \vspace*{0.8cm}
    \begin{center}
      {\LARGE \textbf{Formula Student Monocoque SES Auslegung}}\\[3.2cm]
    
      {\huge \textbf{USER GUIDE}}\\[0.5cm]
      {\Large \textbf{Version \DocVersion}}\\[2.2cm]
      {\Large \textbf{MATLAB Skript}}\\[2.2cm]
      {\large Leitfaden für die Nutzung des Matlabskriptes zur Berechnung \par und Optimierung von SES-Panels}\\[1.8cm]

      {\large Saison 2026 / 2027}\\[1.5cm]

      \textbf{Baugruppe Frame \& Body}\\[0.4cm]
      Laurin Basso\\[1.5cm]

      \vfill
      {\small Dokumentenversion \DocVersion \quad|\quad Stand: 2026}
    \end{center}
  \end{titlepage}
  \clearpage

  % ---------------------------------------------------------------------------------
  % Inhaltsverzeichnis (LaTeX-generiert mit Seitenzahlen)
  % ---------------------------------------------------------------------------------
  \hypersetup{pageanchor=true}
  \pagenumbering{roman}
  \tableofcontents
  \clearpage

  % ---------------------------------------------------------------------------------
  % Tabellenverzeichnis (automatisch generiert & verlinkt)
  % ---------------------------------------------------------------------------------
  \phantomsection
  \addcontentsline{toc}{chapter}{\listtablename}
  \listoftables
  \clearpage

  % ---------------------------------------------------------------------------------
  % Abkürzungsverzeichnis (strukturiert & verlinkt)
  % ---------------------------------------------------------------------------------
  \phantomsection
  \addcontentsline{toc}{chapter}{Abkürzungsverzeichnis}
  \chapter*{Abkürzungsverzeichnis}
\markboth{Abkürzungsverzeichnis}{Abkürzungsverzeichnis}
\vspace{-1em}

\noindent
\begin{tabularx}{\textwidth}{@{}l >{\raggedright\arraybackslash}X@{}}
\phantomsection\label{abk:ASTM}\hyperref[sec:continuous_improvement]{\textbf{ASTM}} & American Society for Testing and Materials (Prüfnorm, z.\,B. ASTM C393 für Sandwich-Biegeversuche) \\[0.4ex]
\phantomsection\label{abk:CFK}\hyperref[sec:systemzweck]{\textbf{CFK}}   & Kohlenstofffaserverstärkter Kunststoff (engl. \textit{CFRP}, Carbon Fiber Reinforced Polymer) \\[0.4ex]
\phantomsection\label{abk:CLPT}\hyperref[sec:systemgrenzen]{\textbf{CLPT}} & Classical Laminated Plate Theory (Klassische Laminattheorie nach Kirchhoff/Mindlin) \\[0.4ex]
\phantomsection\label{abk:CP}\hyperref[sec:fertigungsrichtlinien]{\textbf{CP}}     & Coupling Film (Struktureller Epoxidharz-Klebefilm zur Wabenkernanbindung, z.\,B. CP006) \\[0.4ex]
\phantomsection\label{abk:CSV}\hyperref[sec:materials_db]{\textbf{CSV}}   & Comma-Separated Values (Dateiformat der Werkstoffdatenbank \texttt{materials.csv}) \\[0.4ex]
\phantomsection\label{abk:EI}\hyperref[sec:calc_single_panel]{\textbf{EI}}     & Äquivalente Biegesteifigkeit ($E \cdot I$, Produkt aus Biegemodul und Flächenträgheitsmoment) \\[0.4ex]
\phantomsection\label{abk:FSG}\hyperref[sec:fsg_rules]{\textbf{FSG}}   & Formula Student Germany (Internationales Konstruktions- und Sicherheitsregelwerk) \\[0.4ex]
\phantomsection\label{abk:GUI}\hyperref[sec:sourcetree]{\textbf{GUI}}   & Graphical User Interface (Grafische Benutzeroberfläche, z.\,B. Atlassian Sourcetree) \\[0.4ex]
\phantomsection\label{abk:HM}\hyperref[sec:materials_db]{\textbf{HM}}     & High Modulus (Hochmodul-Kohlenstofffaser mit hoher Steifigkeit, z.\,B. M40J) \\[0.4ex]
\phantomsection\label{abk:IM}\hyperref[sec:materials_db]{\textbf{IM}}     & Intermediate Modulus (Intermediat-Modul-Faser mit hoher Zugfestigkeit, z.\,B. T800S) \\[0.4ex]
\phantomsection\label{abk:L}\hyperref[sec:calibration]{\textbf{L}}       & Ribbon- bzw. Längsrichtung bei Wabenkernen (Richtung höchster Kernschubsteifigkeit) \\[0.4ex]
\phantomsection\label{abk:PBO}\hyperref[sec:materials_db]{\textbf{PBO}}   & Poly(p-phenylen-2,6-benzobisoxazol) (Schlagzähe Hochleistungsfaser Zylon) \\[0.4ex]
\phantomsection\label{abk:PDF}\hyperref[sec:optimize_panel]{\textbf{PDF}}   & Portable Document Format (Druckfertiges Exportformat für Prüf- und Optimierungsberichte) \\[0.4ex]
\phantomsection\label{abk:RAM}\hyperref[sec:ts_cache]{\textbf{RAM}}   & Random-Access Memory (Arbeitsspeicher, vgl. MATLAB-Funktionscache \texttt{clear functions}) \\[0.4ex]
\phantomsection\label{abk:SES}\hyperref[sec:systemzweck]{\textbf{SES}}   & Structural Equivalency Spreadsheet (Reglement-Nachweisdokument der Monocoque-Sicherheit) \\[0.4ex]
\phantomsection\label{abk:UD}\hyperref[tab:quickstart_materials]{\textbf{UD}}     & Unidirektional (Fasergelege mit paralleler Faserorientierung zur Hauptlastaufnahme) \\[0.4ex]
\phantomsection\label{abk:UTS}\hyperref[sec:calc_single_panel]{\textbf{UTS}}   & Ultimate Tensile Strength (Zugfestigkeit bzw. Deckhaut-Bruchspannung im Biegeversuch) \\[0.4ex]
\phantomsection\label{abk:W}\hyperref[sec:calibration]{\textbf{W}}       & Width- bzw. Querrichtung bei Wabenkernen (senkrecht zur Ribbon-Richtung) \\
\end{tabularx}

  \clearpage

  % ---------------------------------------------------------------------------------
  % Hauptteil (Arabische Seitennummerierung)
  % ---------------------------------------------------------------------------------
  \pagenumbering{arabic}

  \chapter{Erste Schritte, Systemgrenzen \& Quick Start} \label{ch:quickstart}

\section{Grundverständnis \& Systemgrenzen} \label{sec:systemgrenzen}

\begin{tcolorbox}[colback=warnbg,colframe=warnred,arc=2mm,title=\textbf{ACHTUNG: Theoretische Simulation vs. Reale Fertigung}]
Die Berechnungen dieses Tools basieren auf analytischen Modellen unter Annahme idealer Materialeigenschaften. 
Die realen Prüfstandsergebnisse weichen in Abhängigkeit von der \textbf{Fertigungsqualität und äußeren Einflüssen} ab.
\end{tcolorbox}

\begin{itemize}
    \item \textbf{Einfluss der Fertigungsqualität}: Faktoren wie der effektive Anpressdruck im Heißpressverfahren, Temperaturzyklen, Faserauslenkung, Porengehalt und Harzaufnahme der offenen Wabenkammern beeinflussen die Steifigkeit und Festigkeit maßgeblich.
    \item \textbf{Erwartete Abweichung}: Eine Abweichung von \textbf{ca. 10\,\%} zwischen Simulation und Prüfstandsergebnis ist bei sorgfältiger, reproduzierbarer Prepreg-Fertigung normal und als Auslegungskorridor zu berücksichtigen.
    \item \textbf{Nichtlinearität der Energieaufnahme}: Während der elastische Anstieg Gradient sehr präzise vorhergesagt werden kann, unterliegt die Energieberechnung bis zur Normdurchbiegung von $12{,}7\,\text{mm}$ größeren Streuungen. 
        Grund hierfür ist der hochgradig nichtlinearität im gesammten Bereich vor und nach Überschreiten der Maximalkraft. Zur fundierten Plausibilisierung der absorbierten Energie wurden diskrete Versagensmodi implementiert (weiterführende Details siehe Technische Dokumentation).   
\end{itemize}
\clearpage










\section{Kurzbeschreibung \& Systemzweck} \label{sec:systemzweck}

Das Tool dient der schnellen, unkomplizierten und regelkonformen Vorauslegung von Faserverbund-Sandwichplatten (\hyperref[abk:CFK]{CFK}-Deckschichten mit Aluminium-Wabenkern oder Hartschaum) für das Formula Student Monocoque.

Es bildet den genormten \textbf{3-Punkt-Biegeversuch} (nach \hyperref[abk:ASTM]{ASTM} C393 / D7250) rechnerisch ab und ermittelt die für das \hyperref[abk:SES]{\textbf{Structural Equivalency Spreadsheet (SES)}} geforderten Strukturkennwerte:
\begin{enumerate}
    \item Äquivalente Biegesteifigkeit (\hyperref[abk:EI]{$EI$}) des Verbunds
    \item Elastisches Biegemodul und Steifigkeitsgradient
    \item Maximale Bruchkraft ($F_{max}$) unter Berücksichtigung Versagensmodi
    \item Energieabsorption ($E_{12.7}$) bis $12{,}7\,\text{mm}$ Stempelweg
    \item und weitere informative Größen
\end{enumerate}










\section{Repository klonen mit Sourcetree} \label{sec:sourcetree}

Um die Versionsverwaltung über Git durch eine übersichtliche und intuitiv bedienbare grafische Benutzeroberfläche (\hyperref[abk:GUI]{GUI}) zu vereinfachen, wird die Verwendung von Atlassian Sourcetree empfohlen:

\begin{enumerate}
    \item \textbf{Insatlation}: https://www.atlassian.com/software/sourcetree
    \item \textbf{Optional GitHub Account}: https://github.com 
    \begin{itemize}
        \item \textbf{Favorisieren}: Favorisiere das Repository um es leichter wieder zu finden.
    \end{itemize}
    \item \textbf{Repository-URL}: https://github.com/Basso-RS/SES-Panel-Berechnung.git 
    \item \textbf{Sourcetree öffnen}: Klicke oben auf die Registerkarte \textbf{„Klonen / Neu“} (\textit{Clone}).
    \item \textbf{Pfade angeben}:
    \begin{itemize}
        \item \textbf{Quellpfad / URL}: Füge die kopierte Git-URL ein.
        \item \textbf{Zielfeld}: Wähle dein lokales Arbeitsverzeichnis.
    \end{itemize}
    \item \textbf{Klonen ausführen}: Klicke auf \textbf{„Klonen“}. Sourcetree lädt alle Dateien, Skripte und Materialdatenbanken herunter.
    \item \textbf{Branch prüfen}: Stelle sicher, dass du auf dem Haupt-Branch (\texttt{main} bzw. \texttt{master}) arbeitest.
\end{enumerate}
\clearpage









\section{Quick Start: In 4 Schritten zur ersten Berechnung} \label{sec:quickstart}

Das gesamte Pfad- und Umgebungsmanagement ist automatisiert! 
\par In vier Schritten führst du deine erste Simulation durch:





\subsection*{Schritt 1: Projekt in MATLAB initialisieren}

Öffne MATLAB und navigiere in das geklonte Projektverzeichnis. \par
Die Projektdatei \texttt{SES\_Panelberechnung.prj} im Verzeichnisbaum  von MATLAB (Links) per Doppelklick oder Rechtsklick ausführen:

\textit{(Hintergrund: Durch das Öffnen des MATLAB-Projekts werden alle erforderlichen Unterordner (\texttt{00\_Scripts}, \texttt{01\_Functions}, \texttt{02\_Data}) automatisch in den MATLAB-Suchpfad aufgenommen.)}





\subsection*{Schritt 2: Panel \& Lagenaufbau prüfen}
Öffne die Datei \texttt{01\_Functions/Scenariodata.m}. \par
Hier ist der zu berechnende Prüfkörper definiert:
\begin{lstlisting}[style=vscodeLightMatlab]
Scenario.panel.b     = 275;   % mm Panel-Breite
Scenario.panel.l     = 400;   % mm Stützweite
Scenario.panel.w     = 12.7;  % mm Max. Durchbiegung
\end{lstlisting}

Layup Informationen für eine Berechnung eines Panels:

\begin{lstlisting}[style=vscodeLightMatlab]
Scenario.name   = '26_16';
Scenario.core_t = 18;        % 18 mm Wabenkern
Scenario.layup  = '45t7/0m/0m/0m/21m/45t7/cp/18core/cp/45t7/-21m/0m/0m/0m/45t7';
\end{lstlisting}
Passe bei Bedarf den Namen, die Kerndicke (\texttt{core\_t}) oder den Layup-String an und speichere die Datei.

Zu beachten ist hier die angegebene Abkürzungskonvention der Material- und Lagen-Tokens:
\begin{itemize}
    \item \textbf{Faserlagen}: Zusammensetzung aus \texttt{[Winkel][Kürzel]}, z.\,B. \texttt{45t7} ($+45^\circ$ T700-Gewebe) oder \texttt{-21m} ($-21^\circ$ \hyperref[abk:UD]{M40J-UD}).
    \item \textbf{Wabenkern}: Deklaration über \texttt{[Dicke]core}, z.\,B. \texttt{18core} für $18\,\text{mm}$ Aluwabe.
    \item \textbf{Klebefilm}: \hyperref[abk:CP]{\texttt{cp}} deklariert den beidseitig am Kern verwendeten Harzfilm (CP006).
\end{itemize}

\begin{table}[htbp]
\centering
\small
\begin{tabular}{llcp{7.5cm}}
\toprule
\textbf{Kürzel} & \textbf{Materialbezeichnung} & \textbf{Typ} & \textbf{Funktion / Eigenschaften} \\
\midrule
\texttt{t7}   & T700 Fabric         & Gewebe & Zähe Sichtlage ($45^\circ$) \& Kernanbindung \\
\texttt{m}    & M40J UD             & UD     & Hauptlastaufnahme ($0^\circ, \pm 21^\circ$, Hochmodul) \\
\texttt{f}    & M40J Fabric         & Gewebe & Hochsteife Decklage (Hochmodul-Gewebe) \\
\texttt{t8}   & T800S UD            & UD     & Hochfeste Zwischenlagen (IM-Faser) \\
\texttt{z}    & ZZ300T              & Gewebe & Zylon/PBO für maximale Energieaufnahme \\
\texttt{a}    & SIGAPREG A W173     & Gewebe & Aramid/Kevlar für Durchschlagschutz \\
\texttt{cp}   & Coupling Film CP006 & Harz   & Klebefilm zwischen Kern und Deckhaut \\
\texttt{core} & PAMG 5056 Aluwabe   & Kern   & Wabenkern (z.\,B. \texttt{18core}) \\
\bottomrule
\end{tabular}
\caption{Materialkürzel und Basisfunktionen im Lagenaufbau.}
\label{tab:quickstart_materials}
\end{table}


\subsection*{Schritt 3: Skript ausführen}

Führe das Berechnungsskript \texttt{Calculate\_Single\_Panel}, 
aus dem Ordner \texttt{00\_Scripts} über RUN im Editor-Tab aus. 

\textit{(Alternativ: Datei \texttt{00\_Scripts/Calculate\_Single\_Panel.m} im MATLAB-Editor öffnen und \texttt{F5} drücken).}





\subsection*{Schritt 4: Ergebnisse einsehen}
Nach Abschluss der Berechnung (Rechenzeit typischerweise $< 0{,}1\,\text{Sekunden}$):
\begin{enumerate}
    \item \textbf{Terminal-Report}: Im MATLAB Command Window erscheint eine detaillierte Zusammenfassung aller Steifigkeits-, Festigkeits- und Kennwerte.
    \item \textbf{Automatischer PDF-Prüfbericht}: Das Skript erzeugt automatisch einen DIN-A4-Prüfbericht und legt diesen mit Zeitstempel im persönlichen \textbf{\texttt{Downloads}}-Verzeichnis ab:
    \begin{quote}
        \texttt{SES\_Report\_Single\_Panel\_26\_20\_<Zeitstempel>.pdf}
    \end{quote}
\end{enumerate}

  \chapter{Erweiterte Konfiguration \& Kalibrierung} \label{ch:konfiguration}

Nachdem im Quick Start (Abschnitt \ref{sec:quickstart}) die grundlegende Geometrie und das Erstellen eines Standard-Laminats in \texttt{Scenariodata.m} vorgestellt wurden, behandelt dieses Kapitel weiterführende Syntax-Regeln, Sonderfunktionen und die physikalischen Kalibrierungsparameter.

\section{Erweiterte Layup-Syntax \& Sonderregeln} \label{sec:layup_syntax}

Über die Standard-Definition hinaus stehen in \texttt{Scenariodata.m} folgende Funktionen und Konventionen zur Verfügung:

\begin{enumerate}
    \item \textbf{Symmetrische Kurzschreibweise (\texttt{sym})}:
    \begin{itemize}
        \item Ist die Unterseite exakt spiegelsymmetrisch zur Oberseite aufgebaut, kann nach dem Kern das Schlüsselwort \texttt{sym} gesetzt werden.
        \item Der Parser spiegelt automatisch alle Lagen vor dem Kern spiegelbildlich auf die Unterseite und kehrt dabei die Vorzeichen von Winkellagen automatisch um (z.\,B. $+21^\circ \to -21^\circ$).
    \end{itemize}
\begin{lstlisting}[style=vscodeLightMatlab]
% Kompakte symmetrische Notation:
Scenario.layup = '45t7/0m/21m/cp/18core/sym';

% Entspricht vollstaendig aufgeloest:
% '45t7/0m/21m/cp/18core/cp/-21m/0m/45t7'
\end{lstlisting}

    \item \textbf{Kernausrichtung \& Beliebige Kerne (\texttt{core}, \texttt{corew}, \texttt{rohacell}, \texttt{nomex})}:
    \begin{itemize}
        \item \textbf{Standard} (\texttt{18core}): Ribbon- bzw. L-Richtung parallel zur Längsachse ($G_L = 550\,\text{MPa}$). Dies ist der Standard für maximale Tragfähigkeit.
        \item \textbf{Querausrichtung} (\texttt{18corew}): Einbau um $90^\circ$ gedreht (W-Richtung). Der Schubmodul sinkt auf $G_W = 230\,\text{MPa}$, was zu spürbar höherer Schubverformung führt.
        \item \textbf{Beliebige Kernmaterialien}: Neben der Aluwabe können beliebige Kerne (z.\,B. \texttt{15rohacell}, \texttt{20nomex}) über \texttt{materials.csv} frei definiert und genutzt werden (siehe Schritt-für-Schritt-Anleitung in Abschnitt \ref{sec:add_materials}).
    \end{itemize}

    \item \textbf{Vorrichtungsspiel \& Setzweg (\texttt{panel.slack})}:
    \begin{itemize}
        \item In \texttt{Scenario.panel.slack = 0.0;} kann ein Setzweg der Prüfaufnahme in $\text{mm}$ hinterlegt werden.
        \item Dieser Versatz verschiebt den Kraftanstieg auf der Wegachse und dient dem exakten Abgleich mit Prüfstandskurven, falls bei Prüfbeginn Vorrichtungsspiel vorliegt.
    \end{itemize}

    \item \textbf{Schichtreihenfolge \& Belastungsrichtung}:
    \begin{itemize}
        \item Die Notation im Layup-String erfolgt strikt von \textbf{oben nach unten}:
        \item \textbf{Vor dem Kern (Oberseite)}: Kontaktzone mit dem Biegestempel $\to$ Druckbeanspruchung im 3-Punkt-Biegeversuch.
        \item \textbf{Nach dem Kern (Unterseite)}: Auflage auf den Stützrollen $\to$ Zugbeanspruchung.
        \item \textit{(Wichtig: Bei asymmetrischen Aufbauten entscheidet diese Orientierung maßgeblich darüber, ob Faserzug- oder Faserdruckversagen die Maximalkraft limitiert.)}
    \end{itemize}
\end{enumerate}

\clearpage

\section{Kalibrierungs- \& Fertigungseinstellungen (Calibration.m)} \label{sec:calibration}

In der Datei \texttt{01\_Functions/Calibration.m} sind die zentralen Kalibrierungswerte für Fertigungseinflüsse, Materialkorrekturen und das Versagensmodell hinterlegt.

\begin{tcolorbox}[colback=warnbg,colframe=warnred,arc=2mm,title=\textbf{ACHTUNG: Zentrale Modellparameter}]
Diese Parameter bilden das physikalische Fundament des gesamten Simulationsmodells. \\
\textbf{Änderungen dürfen nur mit triftigem Grund vorgenommen werden}, da sie direkte Auswirkungen auf alle Berechnungsergebnisse und die Validierung gegen Prüfstandsergebnisse haben!
\end{tcolorbox}

\begin{lstlisting}[style=vscodeLightMatlab]
% Auszug aus 01_Functions/Calibration.m:
Calib.thickness_factor    = 0.95;   % Bauteildicke nach Aushaertung (95% Nenndicke, erhoeht E-Modul um 1/0.95)
Calib.core_dir            = 'L';    % Wabenkern-Richtung ('L' = Ribbon, 'W' = Quer, Fallback)
Calib.include_core_shear  = true;   % Kernschubverzerrung nach ASTM C393 Dreipunktbiegung (Timoshenko)
Calib.fixture_stiffness   = 22218;  % [N/mm] Exakte SES Rig Compliance aus Steel-Tube-Kalibrierversuch
Calib.post_fail_factor    = 0.250;  % Restkraftplateau nach Peak (F_post = F_max * factor)
Calib.failure_criterion   = 'hashin';% 'hashin' oder 'fixed_strain'
Calib.softening_factor    = 0.95;   % Sekanten-Softeningfaktor vor Peak
Calib.indentation_depth   = 0.15;   % Lokale Rolleneindrueckung [mm]
Calib.energy_mode         = 'ses_total';% 'ses_total' (FS-Reglement bis 12.7 mm) | 'to_peak' (Abgleich bis F_max)
\end{lstlisting}

\subsection*{Parameter im Detail für begründete Anpassungen:}

\begin{itemize}
    \item \textbf{Kompaktierungsfaktor} (\texttt{thickness\_factor}): \par
    Berücksichtigt die reale Kompaktierung des Prepregs unter Druck im Heißpressverfahren ($\approx 6\,\text{bar}$). Ein Wert von $0{,}95$ bedeutet, dass das ausgehärtete Laminat $95\,\%$ der theoretischen Nenndicke erreicht. Durch die Verringerung der Schichtdicke bei gleichbleibender Fasermasse steigt der wirksame E-Modul physikalisch um $1 / 0{,}95$.

    \item \textbf{Prüfstands- \& Vorrichtungssteifigkeit (Rig Compliance)} (\texttt{fixture\_stiffness}): \par
    Reale Prüfmaschinen und Biegeaufnahmen verformen sich unter Prüfkraft elastisch. Im Berechnungsmodell wird diese Nachgiebigkeit über eine Reihenschaltung der Federsteifigkeiten als "`virtueller Prüfstand"' abgebildet:
    \begin{equation}
        \frac{1}{K_{\text{ges}}} = \frac{1}{K_{\text{Biegung}}} + \frac{1}{K_{\text{Schub}}} + \frac{1}{K_{\text{rig}}} = \frac{L^3}{48 \, EI} + \frac{L}{4 \, S} + \frac{1}{K_{\text{rig}}}
    \end{equation}
    Der Wert $K_{\text{rig}} = 22\,218\,\text{N/mm}$ stammt direkt aus dem offiziellen Formula Student SES Baseline Steel Tube Test (T3.2) mit zwei kalibrierten Stahlrohren auf $400\,\text{mm}$ Stützweite. Dadurch bleiben alle Messdaten in der Excel-Tabelle $100\,\%$ unverfälschte Rohdaten, während die Simulation die reale Maschinenelastizität physikalisch exakt einrechnet.

    \item \textbf{Kernschubanteil nach Timoshenko} (\texttt{include\_core\_shear}): \par
    Berücksichtigt neben der reinen Deckhautbiegung ($L^3 / (48 EI)$) die transversale Kernschubverzerrung nach \abk{ASTM} C393 ($L / (4 S)$). Bei Sandwichplatten mit Wabenkern kann der Schubanteil je nach Kerndicke und Faseraufbau $15\,\%$ bis $30\,\%$ der elastischen Gesamtdurchbiegung ausmachen.

    \item \textbf{Kernausrichtung} (\texttt{core\_dir}): \par
    Hauptausrichtung des Wabenkerns.
    \begin{itemize}
        \item \hyperref[abk:L]{\texttt{'L'}}: Längs-/Ribbon-Richtung ($G_{core} = 550\,\text{MPa}$) für maximale Tragfähigkeit in Längsrichtung.
        \item \hyperref[abk:W]{\texttt{'W'}}: Querrichtung ($G_{core} = 230\,\text{MPa}$). Führt zu deutlich höherer Schubverformung und geringerem Steifigkeitsgradienten.
    \end{itemize}

    \item \textbf{Versagenskriterium} (\texttt{failure\_criterion}): \par
    Bruchhypothese für die Deckschichten.
    \begin{itemize}
        \item \texttt{'hashin'}: Differenziert zwischen Faserzug, Faserdruck, Matrixzug und Matrixdruck unter Spannungsinteraktion.
        \item \texttt{'fixed\_strain'}: Einfacher Dehnungsnachweis bei Erreichen einer vorgegebenen Grenzdehnung.
    \end{itemize}

    \item \textbf{Nachbruchfaktor} (\texttt{post\_fail\_factor}): \par
    Definiert das Restkraft-Niveau nach dem ersten Versagens-Peak ($F_{post} = F_{max} \cdot \text{factor}$). Dieser Wert wurde experimentell über alle Prüfstandsdaten der Saisons 2024--2027 kalibriert ($0{,}250 = 25\,\%$ Resttragfähigkeit).

    \item \textbf{Energiemodus} (\texttt{energy\_mode}): \par
    Definiert das Integrationsintervall für die absorbierte Verformungsenergie:
    \begin{itemize}
        \item \texttt{'to\_peak'}: Berechnet die Energieaufnahme bis zum ersten Faserbruch bei $F_{max}$. Wird genutzt, um die Berechnung 1:1 mit den Messwerten aus der Prüfstands-Excel (\texttt{02\_Data/\allowbreak Pruefstand\_\allowbreak Rohdaten\_\allowbreak Benchmark.xlsx}) bis zum Peak abzugleichen.
        \item \texttt{'ses\_total'}: Vollständige Energieintegration über den gesamten Weg von $0$ bis $12{,}7\,\text{mm}$ (inklusive Wabenstauchung nach dem Bruch). Wird für den offiziellen Formula Student \abk{SES}-Nachweis zwingend gefordert.
    \end{itemize}
\end{itemize}

  \chapter{Die Haupt-Workflows} \label{ch:workflows}

Das Tool bietet zwei primäre Arbeitsabläufe für die Panel-Auslegung: die schnelle Einzelsimulation eines Prüfkörpers und die automatisierte Optimierung zur Gewichtsminimierung unter Einhaltung aller SES- und FSG-Vorgaben.

\section{Workflow 1: Einzelsimulation eines Panels} \label{sec:calc_single_panel}

Dieser Workflow dient der schnellen Berechnung und Überprüfung eines fest definierten Laminataufbaus (z.\,B. zur Überprüfung einer lagenaufbau Idee oder eine Berechnung eines bereits Physisch getesteten Panels).

\subsection*{Ablauf \& Durchführung}
\begin{enumerate}
    \item \textbf{Eingaben prüfen}: In der Datei \texttt{01\_Functions/Scenariodata.m} Geometrie und Lagenaufbau festlegen und speichern (siehe Abschnitt \ref{sec:quickstart}).
    \item \textbf{Skript starten}: Die Datei \texttt{00\_Scripts/Calculate\_Single\_Panel.m} im MATLAB-Editor öffnen und mit \texttt{F5} (oder Klick auf \textbf{Run}) ausführen.
\end{enumerate}

\subsection*{Auswertung des Terminal-Reports}
Die Berechnung wird in wenigen Millisekunden abgeschlossen. Im MATLAB Command Window wird ein detaillierter, in fünf Abschnitte gegliederter Prüfbericht ausgegeben:

\begin{itemize}
    \item \textbf{Geometrie}: Dicken der oberen und unteren Deckhaut, Kerndicke mit Orientierung (Ribbon \texttt{'L'} oder Quer \texttt{'W'}) sowie die Gesamtsandwichdicke.
    \item \textbf{Gewicht \& Flächenmasse}: Gesamtmasse des $400 \times 275\,\text{mm}$ Prüfkörpers in Gramm, gesamtes Flächengewicht in $\text{g/m}^2$ sowie die prozentualen Massenanteile der CFK-Deckhäute, des Wabenkerns und des Klebefilms (\hyperref[abk:CP]{\texttt{cp}}).
    \item \textbf{Steifigkeit (Stiffness)}: Elastizitätsmodul der Deckschichten ($E$), äquivalente Biegesteifigkeit (\hyperref[abk:EI]{$EI$}), axiales Flächenträgheitsmoment ($I$) und der maßgebliche SES-Prüfstandsgradient in $\text{N/mm}$.
    \item \textbf{Festigkeit \& Versagen (Strength \& Failure)}: Deckhaut-Bruchspannung (\hyperref[abk:UTS]{$UTS$}), Maximalkraft ($F_{max}$), Versagensweg bei Maximalkraft ($w_{max}$) sowie der identifizierte primäre Versagensmodus. \par
    \textit{Wichtiger Hinweis zur Modellgenauigkeit}: Je nach auftretendem Versagensmodus unterscheidet sich die Linearität der Kraft-Weg-Kennlinie maßgeblich, wodurch die berechnete Energieaufnahme mehr oder weniger vertrauenswürdig ist:
    \begin{itemize}
        \item \textit{Deckhautbruch (UTS)}: Faserzug- oder Faserdruckversagen der Deckschichten. Durch das überwiegend linear-elastische Verhalten der Fasern bis zum Bruchpunkt weist die Energieberechnung hier die \textbf{höchste Linearität und Vertrauenswürdigkeit} auf.
        \item \textit{Kernschubbruch (Core Shear)}: Schubversagen der Aluminiumwaben bei hoher Querkraft. Plastische Verformungen der Zellwände vor dem Peak mindern die Linearität, wodurch die berechnete Energie moderaten Streuungen unterliegt.
        \item \textit{Deckhautknicken (Wrinkling)}: Lokales Stabilitätsversagen (Beulen) der Deckhaut über den Wabenkammern. Dieses plötzliche, geometrisch nichtlineare Versagen reagiert extrem sensibel auf Fertigungseinflüsse, weshalb der Energieberechnung hier \textbf{am wenigsten vertraut} werden kann.
    \end{itemize}
    \item \textbf{Energieaufnahme (Energy Absorption)}: Absorbierte Verformungsenergie in Joule ($J$) im gewählten Modus (\texttt{'to\_peak'} für den direkten Excel abgleich oder \texttt{'ses\_total'} für den \hyperref[abk:SES]{SES}-Regelwerksnachweis).
\end{itemize}

\subsection*{Automatischer PDF-Prüfbericht}
Das Skript erzeugt automatisch einen DIN-A4-Prüfbericht im \hyperref[abk:PDF]{PDF}-Format und legt diesen mit Zeitstempel im persönlichen Downloads-Verzeichnis ab:
\begin{quote}
    \texttt{SES\_Report\_Single\_Panel\_<Szenarioname>\_<Zeitstempel>.pdf}
\end{quote}
Der Bericht fasst alle Kennwerte übersichtlich zusammen, eignet sich direkt zur Dokumentation und späterem abgleich der realen Prüfstandsergebnisse.

\clearpage

\section{Workflow 2: Automatisierte Laminat-Optimierung} \label{sec:optimize_panel}

Das Skript \texttt{00\_Scripts/Optimize\_Panel.m} führt eine mehrkriterielle Optimierung durch. Ziel ist es, aus allen regelkonformen Kombinationen von Fasermaterialien, Lagenanzahlen und Orientierungswinkeln das leichteste Sandwich-Laminat zu ermitteln, das alle geforderten Strukturkennwerte erfüllt.

\subsection*{Zielwerte definieren (Targets)}
Die Mindestanforderungen an das optimierte Panel werden direkt im Konfigurationsabschnitt von \texttt{Optimize\_Panel.m} hinterlegt:

\begin{lstlisting}[style=vscodeLightMatlab]
Targets = struct();
Targets.safety_factor = 1.0;    % Sicherheitsfaktor auf Zielwerte (z.B. 1.0 = keine Reserve, 1.1 = +10%)
Targets.Gradient      = 3000;   % Mindest-Gradient [N/mm]
Targets.Energy        = 60;     % Mindest-Energieaufnahme [J]
Targets.EI            = 0;      % Mindest-Biegesteifigkeit [N*mm^2]
Targets.UTS           = 0;      % Mindest-Deckhautfestigkeit [MPa]
Targets.E             = 0;      % Mindest-E-Modul Deckhaut [GPa]
Targets.F_max         = 0;      % Mindest-Bruchkraft [N]
Targets.t_skin_max    = 0;      % Max. Deckhautdicke [mm]
\end{lstlisting}

\begin{itemize}
    \item \textbf{Sicherheitsfaktor (\texttt{safety\_factor})}: Skaliert alle aktivierten Mindest-Targets intern mit diesem Faktor (z.\,B. $1{,}10$ für $+10\,\%$ Reserve: $3000 \cdot 1{,}10 = 3300\,\text{N/mm}$).
    \item \textbf{Inaktiv-Schalter ($0$)}: Jeder Zielwert mit dem Wert $0$ wird vom Optimierungsalgorithmus ignoriert und stellt keine Restriktion dar.
\end{itemize}

\subsection*{Fertigungs- \& Regelwerksgrenzen (OptSettings)}
Zur Einhaltung von Fertigungsstandards und FSG-Vorgaben stehen folgende Optionen zur Verfügung:

\begin{lstlisting}[style=vscodeLightMatlab]
OptSettings.min_plies             = 4;           % Mindestanzahl Lagen je Deckhaut
OptSettings.max_plies             = 8;           % Maximalanzahl Lagen je Deckhaut
OptSettings.angle_mode            = 'simple';    % 'simple' ([0, +/-21]) oder 'advanced'
OptSettings.material_pool         = 'all';       % 'all' oder 'standard'
OptSettings.sichtlage_active      = true;        % FSG-Sichtlage zwingend aktiv
OptSettings.sichtlage_material    = {'t7', 'f'}; % Erlaubte Sichtlagengewebe
OptSettings.sichtlage_angle       = 45;          % Sichtlagenwinkel in Grad
OptSettings.core_interface_fabric = true;        % Gewebelage zur Kernanbindung
\end{lstlisting}

\begin{itemize}
    \item \textbf{Lagenbegrenzung} (\texttt{min\_plies} / \texttt{max\_plies}): Definiert den Suchbereich für die Einzellagen pro Deckhaut (Standard: 4 bis 8 Lagen).
    \item \textbf{Winkelmodus} (\texttt{angle\_mode}): \texttt{'simple'} nutzt die bewährten Standardwinkel ($0^\circ, \pm 21^\circ$); \texttt{'advanced'} erweitert den Suchraum um $\pm 30^\circ, \pm 45^\circ, \pm 60^\circ$.
    \item \textbf{Sichtlage} (\texttt{sichtlage\_active}): Garantiert die Einhaltung der FSG-Regel T\,3.5.7 (äußerste Lage muss zwingend einlagiges Gewebe sein, Standard: $+45^\circ$ \texttt{t7} oder \texttt{f}).
    \item \textbf{Kernanbindung} (\texttt{core\_interface\_fabric}): Erzwingt eine Gewebelage direkt am Wabenkern, um Delaminationsrisiken vorzubeugen.
\end{itemize}

\subsection*{Auswertung der Optimierungsergebnisse}
Nach Ausführung von \texttt{Optimize\_Panel.m} werden drei Ergebnisformate bereitgestellt:
\begin{enumerate}
    \item \textbf{Vorher/Nachher-Vergleich im Terminal}: Direkte Gegenüberstellung der Baseline (aus \texttt{Scenariodata.m}) gegen das Optimum mit genauer Massenersparnis in Gramm und Prozent.
    \item \textbf{TOP-5 Rangliste}: Auflistung der 5 leichtesten regelkonformen Laminate inklusive Halbschalen-Layup-String zur direkten Übernahme.
    \item \textbf{Kraft-Weg-Vergleichsplot}: Visuelle Gegenüberstellung der Arbeitskurven von Baseline (Rot) und Optimum (Grün) mit Markierung von $F_{max}$ und der geforderten Gradienten-Grenzlinie.
    \item \textbf{PDF-Optimierungsbericht}: Automatischer Export der gesamten Zusammenfassung als PDF in den \textbf{\texttt{Downloads}}-Ordner.
\end{enumerate}

  \chapter{Troubleshooting \& Fehlerbehebung} \label{ch:troubleshooting}

Dieses Kapitel bietet schnelle Hilfe bei typischen Eingabefehlern, Fehlermeldungen in MATLAB und unplausiblen Berechnungsergebnissen.

\section{Syntax- \& Eingabefehler im Layup} \label{sec:ts_syntax}

Der Parser (\texttt{parse\_layup.m}) zerlegt den Lagenaufbau-String entlang der Trennzeichen \texttt{'/'}. Weicht die Syntax von der Konvention ab, reagiert das System mit Fehlermeldungen oder unerwarteten Modellen:

\begin{itemize}
    \item \textbf{Unbekanntes Materialkürzel}:
    \begin{quote}
        \texttt{Error: Material xyz nicht in materials.csv gefunden.}
    \end{quote}
    \textit{Ursache}: Ein Token wurde falsch getippt oder die Reihenfolge von Winkel und Material wurde vertauscht (z.\,B. \texttt{m0} statt \texttt{0m} oder \texttt{f45} statt \texttt{45f}). Der Parser erwartet immer zuerst den Winkel (inkl. Vorzeichen), direkt gefolgt vom Kürzel.
    
    \item \textbf{Falsche Vorzeichenplatzierung}:
    Ein Token wie \texttt{45-m} wird als Winkel $45^\circ$ und Material \texttt{'-m'} interpretiert. Das Vorzeichen muss zwingend ganz vorne stehen (z.\,B. \texttt{-45m} oder \texttt{-21m}).

    \item \textbf{Vergessener Kern-Token (\texttt{core})}:
    Wird kein Kern (z.\,B. \texttt{18core}) angegeben, interpretiert das Tool den Aufbau als monolithisches Voll-CFK-Laminat ohne Sandwich-Kern. Die Kernhöhe wird zu \texttt{NaN} gesetzt und die Kernschubberechnung entfällt.

    \item \textbf{Fehlender Klebefilm (\texttt{cp})}:
    Wird \texttt{cp} vor und nach dem Kern weggelassen, bricht die Simulation zwar nicht ab, aber die berechnete Panelmasse ist ca. $100\,\text{g/m}^2$ pro Deckhaut zu niedrig.

    \item \textbf{Fehlerhafte Symmetrie-Syntax (\texttt{sym})}:
    Der Token \texttt{sym} darf ausschließlich als letztes Element nach dem Kern stehen (z.\,B. \texttt{.../18core/cp/sym}). Werden nach \texttt{sym} weitere Lagen notiert, wird das Schlüsselwort als Material-Token fehlinterpretiert.
\end{itemize}

\section{Materialdaten \& Datenbankprobleme} \label{sec:ts_cache}

Die Materialwerte werden zentral über die Datei \texttt{02\_Data/materials.csv} verwaltet. Beim Bearbeiten oder Hinzufügen von Datensätzen sind folgende Punkte zu beachten:

\begin{itemize}
    \item \textbf{Werteänderung in \hyperref[abk:CSV]{CSV} wird in MATLAB nicht übernommen}:
    Die Funktion \texttt{load\_materials.m} nutzt aus Performancegründen einen internen \hyperref[abk:RAM]{RAM}-Cache. Wird die CSV-Datei extern in Excel oder einem Editor geändert, greift MATLAB weiterhin auf die im Speicher liegenden Altwerte zu. \par
    \textit{Lösung}: Vor einem erneuten Rechenlauf den Speicher über das Command Window leeren:
    \begin{lstlisting}[style=vscodeLightMatlab]
    clear functions; % Leert den Material-Cache
    \end{lstlisting}
    Alternativ kann das Neuladen mit \texttt{load\_materials([], true)} erzwungen werden.

    \item \textbf{Falsches CSV-Trennzeichen}:
    Die Datei \texttt{materials.csv} muss zwingend mit Semikolon (\texttt{;}) als Spaltentrennzeichen gespeichert werden. Werden Kommata oder Tabulatoren verwendet, schlägt das Einlesen fehl (\texttt{detectImportOptions}).

    \item \textbf{Fehlende Festigkeitskennwerte in CSV}: \par
    Sind Kennwerte wie $X_t, X_c, Y_t, Y_c$ oder $S_{12}$ in der CSV-Tabelle leer oder $\le 0$, setzt der Parser automatisch einen Schutzwert von $10^6\,\text{MPa}$. In diesem Fall meldet das Tool keinen Faserbruch für diese Lage.
\end{itemize}

\section{Unplausible Simulationsergebnisse} \label{sec:ts_ergebnisse}

Weichen die Ergebnisse drastisch von Erfahrungswerten oder Prüfstandsdaten ab, liegt meist eine Fehleinstellung in den Parametern vor:

\begin{itemize}
    \item \textbf{Biegesteifigkeit oder Gradient deutlich zu niedrig}:
    \begin{itemize}
        \item \textit{Kernorientierung}: Steht der Kern auf Querausrichtung (\texttt{'W'}), halbiert sich der Schubmodul von $550\,\text{MPa}$ auf $230\,\text{MPa}$. Dadurch wachsen die Schubverformungen massiv an.
        \item \textit{Kompaktierung}: Wurde \texttt{thickness\_factor} in \texttt{Calibration.m} versehentlich verstellt? Der Standardwert liegt bei $0{,}95$.
    \end{itemize}

    \item \textbf{Kraft-Weg-Kurve ist horizontal verschoben}:
    Liegt die simulierte Kurve parallel versetzt zur realen Prüfstandskurve, ist in \texttt{Scenariodata.m} ein Vorrichtungsspiel eingetragen (\texttt{Scenario.panel.slack}). Für theoretische Panel-Auslegungen gehört der Wert auf $0{,}0\,\text{mm}$.

    \item \textbf{Berechnete Energie passt nicht zum Prüfstandsversuch}:
    \begin{itemize}
        \item \textit{Energiemodus prüfen}: In \texttt{Calibration.m} prüfen, ob \texttt{'to\_peak'} (Abgleich bis $F_{max}$) oder \texttt{'ses\_total'} ($0$ bis $12{,}7\,\text{mm}$ inkl. Restplateau) aktiv ist.
        \item \textit{Versagensmodus beachten}: Tritt Deckhautknicken (Wrinkling) auf, ist der Realversuch stark von geometrischen Fertigungsungenauigkeiten geprägt. Die theoretische Energieberechnung besitzt hier konstruktionsbedingt die größte Unsicherheit.
    \end{itemize}
\end{itemize}

\section{Probleme bei der Laminat-Optimierung} \label{sec:ts_performance}

Beim Ausführen des Skripts \texttt{00\_Scripts/Optimize\_Panel.m} können folgende Konstellationen auftreten:

\begin{itemize}
    \item \textbf{Meldung: \texttt{KEINE zulässige Lösung gefunden, die alle Targets erfüllt!}}:
    \begin{itemize}
        \item \textit{Sicherheitsfaktor zu hoch}: Bei aktivem Sicherheitsfaktor (z.\,B. \texttt{safety\_factor = 1.25}) steigen die Anforderungen um $25\,\%$.
        \item \textit{Lagenfenster zu eng}: Wenn \texttt{max\_plies = 6} gewählt ist, der geforderte Gradient aber mindestens 8 Lagen verlangt, existiert im gewählten Suchraum keine physikalisch machbare Lösung. \texttt{max\_plies} schrittweise anheben.
        \item \textit{Widersprüchliche Randbedingungen}: Ein hoher Gradient in Kombination mit einer extrem strengen Dickenbegrenzung (\texttt{t\_skin\_max}) schließt alle tragfähigen Laminate aus.
    \end{itemize}

    \item \textbf{Lange Rechenzeit bei der Optimierung}: \par
    Steht \texttt{angle\_mode} auf \texttt{'advanced'}, untersucht der Algorithmus sämtliche Winkelkombinationen aus $[21^\circ, 30^\circ, 45^\circ, 60^\circ]$. Dadurch wächst die Anzahl der Kandidaten exponentiell an. Für schnelle Iterationen immer \texttt{'simple'} nutzen.
\end{itemize}

\section{Dateisystem, Pfade \& PDF-Export} \label{sec:ts_export}

\begin{itemize}
    \item \textbf{\texttt{Unrecognized function or variable}}:
    MATLAB findet die Funktionen nicht. \par
    \textit{Lösung}: Vor Arbeitsbeginn immer die Projektdatei \texttt{SES\_Panelberechnung.prj} im Hauptverzeichnis doppelklicken. Dadurch werden alle Unterordner (\texttt{00\_Scripts}, \texttt{01\_Functions}, \texttt{02\_Data}) automatisch dem MATLAB-Suchpfad hinzugefügt.

    \item \textbf{PDF-Export schlägt fehl oder Datei wird nicht aktualisiert}:
    Die Funktion \texttt{export\_pdf\_report.m} speichert Berichte standardmäßig im Ordner \texttt{Downloads}. Ist ein bereits erzeugter PDF-Bericht in einem externen PDF-Viewer (z.\,B. Adobe Acrobat) geöffnet, sperrt Windows die Datei gegen Überschreiben. \par
    \textit{Lösung}: PDF-Viewer schließen oder das Dokument im Viewer freigeben.
\end{itemize}

  \chapter{Sonderfunktionen \& Fortgeschrittene Workflows} \label{ch:advanced}

Dieses Kapitel behandelt erweiterte Werkzeuge zur empirischen Modellvalidierung, den kontinuierlichen Datenabgleich mit Prüfstandsversuchen sowie fundamentale Fertigungsrichtlinien für den physischen Bau von SES-Testpanels.

\section{Benchmark-Suite zur Modellvalidierung} \label{sec:benchmark}

Das Skript \texttt{Benchmark\_Panels.m} (bzw. die Kernfunktion \texttt{run\_benchmark.m}) dient der kontinuierlichen Überprüfung der Modellgüte. Es führt einen vollautomatisierten Abgleich zwischen den theoretischen Simulationsergebnissen und den realen Messwerten aller gefertigten Formula-Student-Panels durch.

\subsection*{Architektur \& Live-Rohdatenauswertung}
Im Gegensatz zu statischen Tabellenwerten arbeitet die Benchmark-Suite direkt auf den realen Prüfmaschinen-Messdaten:
\begin{itemize}
    \item \textbf{Single Source of Truth}: Sämtliche Panel-Definitionen und Kraft-Weg-Messkurven werden zentral in der Datei \texttt{02\_Data/\allowbreak Pruefstand\_Rohdaten\_Benchmark.xlsx} gepflegt.
    \item \textbf{Automatisierte Messkurven-Analyse}: Die Funktion \texttt{evaluate\_test\_curve.m} liest die Rohdaten (\texttt{Traversenweg\_mm} und \texttt{Kraft\_N}) für jedes Panel ein. Sie bereinigt Vorlast- und Setzeffekte automatisch, ermittelt die Maximalkraft $F_{\max}$, berechnet über eine lineare Regression (im Bereich $15\,\%$ bis $45\,\% \, F_{\max}$) den realen Steifigkeitsgradienten und integriert die absorbierte Energie exakt bis zum \abk{SES}-Normstempelweg von $12{,}70\,\text{mm}$.
    \item \textbf{Identische Berechnungskette}: Für jedes Panel ruft das Benchmark-Tool exakt die modulare Standard-Simulationskette des Projekts auf (\texttt{Scenariodata}, \allowbreak \texttt{Calibration}, \allowbreak \texttt{Materialdata}, \allowbreak \texttt{EI\_calc}, \allowbreak \texttt{Shearstiff\_calc}, \allowbreak \texttt{Deflection}, \allowbreak \texttt{Energy\_calc}). Es finden keinerlei isolierte Sonderberechnungen statt.
    \item \textbf{Duale Deckhautdicken-Erfassung für asymmetrische Panels}: Um asymmetrische Monocoque-Panels (z.\,B. unterschiedliche Lagendicken auf Druck- und Zugseite gemäß \abk{FSG} T 3.5.4) exakt abzugleichen, erfasst der Benchmark getrennte Messwerte für Obergurt und Untergurt (\texttt{Deckhautdicke\_oben\_mm} und \texttt{Deckhautdicke\_unten\_mm}). Die Suite vergleicht beide Seiten individuell mit der Laminattheorie und schaltet die Tabellenausgabe bei asymmetrischen Aufbauten automatisch in die getrennte Darstellung \texttt{oben / unten}.
\end{itemize}

\subsection*{Aktiver Datenbestand (28 Panels)}
Der Benchmark filtert und validiert aktuell 28 aussagekräftige Fertigungspanels aus drei Saisons:
\begin{itemize}
    \item \textbf{Saison 2024 (17 Panels)}: Fertigungspanels \texttt{24\_02} bis \texttt{24\_18} (ohne die irregulären Proben \texttt{24\_19} und \texttt{24\_20}).
    \item \textbf{Saison 2026 (7 Panels)}: Aktuelle Monocoque-Panels \texttt{26\_12} bis \texttt{26\_21} (unter Ausschluss von Frühausfällen durch Fertigungsfehler wie Delamination bei \texttt{26\_17} und Sonderwerkstoffen bei \texttt{26\_19}).
    \item \textbf{Saison 2027 (4 Panels)}: Neueste Prüfpanels \texttt{27\_01} bis \texttt{27\_04} inklusive der Panels mit $20\,\text{mm}$ Wabenkern.
\end{itemize}

\subsection*{Prioritäre Zielgrößen \& Korrelationsgüte}
Für das \abk{SES} sind zwei Kennwerte von überragender Bedeutung:
\begin{enumerate}
    \item \textbf{Steifigkeitsgradient ($\text{Gradient}$ in $\text{N/mm}$)}: Maßgeblich für die Einhaltung der Mindeststeifigkeit des Cockpits und der Seitenwände.
    \item \textbf{Energieaufnahme ($\text{Energy}$ in $\text{J}$)}: Maßgeblich für den Nachweis der Crash- und Verformungsenergie bis $12{,}7\,\text{mm}$.
\end{enumerate}
Zusätzlich dient der Dickenabgleich der Fertigungskontrolle: Bei symmetrischen Aufbauten erfolgt der Vergleich kompakt einzeilig ($t_{\text{test}}$ vs. $t_{\text{calc}}$), während bei asymmetrischen Panels beide Deckhäute getrennt als $t_{\text{top}}/t_{\text{bot}}$ mit individuellen prozentualen Abweichungen ausgewiesen werden (z.\,B. \texttt{0.65/0.65} vs. \texttt{0.84/0.64\,mm}). Die Benchmark-Statistik dokumentiert dies als \texttt{MAE Dicke (o/u)}.

Dank der Reihenschaltung von Plattenbiegung, Wabenkern-Schubverzerrung nach \abk{ASTM} C393 und der Prüfstands-Compliance ($K_{\text{rig}} = 22\,218\,\text{N/mm}$ aus dem SES Baseline Steel Tube Test T3.2) sank die mittlere absolute Abweichung (MAE) über alle 28 Panels auf ca. $12{,}7\,\%$ für den Gradienten und unter $10\,\%$ bei den aktuellen Monocoque-Panels der Saison 2026.

Nach jedem Gesamtlauf erzeugt das Skript automatisch einen druckfertigen DIN-A4-Prüfbericht im \abk{PDF}-Format im Ordner \textbf{\texttt{Downloads}}:
\begin{quote}
    \texttt{SES\_Report\_Benchmark\_Validation\_<Zeitstempel>.pdf}
\end{quote}

\section{Messdaten-Abgleich \& Continuous Improvement} \label{sec:continuous_improvement}

Das Berechnungsmodell ist kein statisches Tool, sondern ein lernendes System, das kontinuierlich an realen Biegeversuchen geschärft wird.

\begin{tcolorbox}[colback=warnbg,colframe=warnred,arc=2mm,title=\textbf{Wichtiger Leitsatz: Pflicht zur lückenlosen Panel-Aufnahme}]
\textbf{Jedes neu gefertigte und im 3-Punkt-Biegeversuch geprüfte SES-Panel MUSS zwingend in die Benchmark-Datenbasis aufgenommen werden!} \par \vspace{0.5em}
\end{tcolorbox}

\subsection*{Workflow: Neues Prüfpanel in 4 Schritten einpflegen}
Sobald ein Prüfkörper auf der Universalprüfmaschine getestet und die Kraft-Weg-Messreihe exportiert wurde, erfolgt die Aufnahme in vier standardisierten Schritten:

\begin{enumerate}
    \item \textbf{Panel-Stammdaten eintragen}: \par
    Öffne die Excel-Datei \texttt{02\_Data/\allowbreak Pruefstand\_Rohdaten\_Benchmark.xlsx} und navigiere zum ersten Tabellenblatt \texttt{Panels}. Füge am Ende der Tabelle eine neue Zeile mit den Parametern des Prüflings ein:
    \begin{itemize}
        \item \texttt{Panel\_ID}: Eindeutiges Saisons-Kürzel (z.\,B. \texttt{27\_05}).
        \item \texttt{Layup}: Der vollständige, syntaktisch korrekte Lagenaufbaustring (z.\,B. \texttt{45t7/\allowbreak 0m/\allowbreak 0m/\allowbreak 21m/\allowbreak 45t7/\allowbreak cp/\allowbreak 18core/\allowbreak sym}).
        \item \texttt{Kerndicke\_mm}: Nominale Kernhöhe in $\text{mm}$ (z.\,B. $18$ oder $20$).
        \item \texttt{Deckhautdicke\_oben\_mm}: Real mit der Mikrometer-Schraube gemessene Dicke der oberen Deckhaut (Druckseite / Stempelkontakt) nach der Heißpressung.
        \item \texttt{Deckhautdicke\_unten\_mm}: Real gemessene Dicke der unteren Deckhaut (Zugseite / Auflage auf den Stützrollen).
    \end{itemize}
    \textit{(Hinweis zu asymmetrischen Panels: Bei asymmetrischen Lagenaufbauten müssen beide Deckhäute getrennt gemessen und eingetragen werden. Bei symmetrischen Laminaten sind beide Werte identisch).}

    \item \textbf{Rohdaten-Tabellenblatt anlegen}: \par
    Erstelle in derselben Arbeitsmappe ein neues Tabellenblatt, dessen Name exakt der \texttt{Panel\_ID} entspricht (z.\,B. Blattname \texttt{27\_05}). \\
    Füge die unformatierten Messwerte der Prüfmaschine spaltenweise ein:
    \begin{itemize}
        \item Spalte A (\texttt{Traversenweg\_mm}): Stempelweg bzw. Durchbiegung $w$ in $\text{mm}$.
        \item Spalte B (\texttt{Kraft\_N}): Gemessene Prüfkraft $F$ in $\text{N}$.
        \item \textit{(Optional Spalte C \texttt{Zeit\_s}: Zeitstempel der Messung).}
    \end{itemize}

    \item \textbf{Benchmark mit Neuladen ausführen}: \par
    Starte den Benchmark in MATLAB mit erzwungenem Cache-Reload, um die neuen Excel-Daten einzulesen:
    \begin{lstlisting}[style=vscodeLightMatlab]
    % Im MATLAB Command Window:
    run_benchmark('reload');
    % Oder alternativ ueber das Gesamtskript:
    Benchmark_Panels;
    \end{lstlisting}

    \item \textbf{Ergebnisse prüfen \& Modell abgleichen}: \par
    Im Terminal-Report und im generierten \abk{PDF}-Bericht die relative Abweichung für Gradient und Energie prüfen:
    \begin{itemize}
        \item \textbf{Normalbereich ($\le 10\,\%$ Abweichung)}: Auslegung und Fertigungsqualität sind bestätigt.
        \item \textbf{Signifikante Abweichung ($> 15\,\%$)}: Prüfen, ob ein unerwarteter Versagensmodus vorlag (z.\,B. Schälung der Deckhaut durch fehlenden Klebefilm \abk{CP} oder Quereinbau des Wabenkerns \abk{W}). Bei systematischen Abweichungen über mehrere Proben können die Fertigungskorrekturfaktoren in \texttt{01\_Functions/Calibration.m} nachjustiert werden.
    \end{itemize}
\end{enumerate}

\section{Fertigungsrichtlinien für den Panelbau} \label{sec:fertigungsrichtlinien}

Aus den Benchmark-Ergebnissen und Schadensanalysen realer Prüfstandstests leiten sich fundamentale Richtlinien für den Panelbau ab:

\begin{itemize}
    \item \textbf{Beidseitiger Klebefilm (\hyperref[abk:CP]{\texttt{cp}})}: \par
    Zwischen Wabenkern und CFK-Deckhaut muss zwingend auf beiden Seiten eine Schicht Epoxidharz-Klebefilm ($\approx 150\,\text{g/m}^2$) eingebracht werden. Ein Weglassen oder unzureichender Anpressdruck führt zu vorzeitiger Deckhaut-Delamination bei Bruchteilen der Auslegungslast.

    \item \textbf{Gewebelage als Kernanbindung}: \par
    Direkt auf den Klebefilm sollte immer mindestens eine Gewebelage (z.\,B. \texttt{45t7} oder \texttt{45f}) gelegt werden. Reine $0^\circ$-UD-Lagen am Kern neigen unter Querkraft zu Schälbeanspruchungen und lokaler Faserabplatzung an den Wabenstegen.

    \item \textbf{Heißpress-Aushärtung}: \par
    Um den im Berechnungsmodell hinterlegten Kompaktierungsfaktor von $0{,}95$ zu erreichen, müssen die Panels in einer Heißpresse unter definiertem Druck ($\approx 6\,\text{bar}$) gehärtet werden. Eine reine Vakuuminfusion oder Vakuumsackhärtung im Ofen erreicht nicht die erforderliche Faservolumenfraktion und führt zu weicheren Deckhäuten.

    \item \textbf{Wabenkern-Orientierung einhalten}: \par
    Vor dem Zuschnitt ist die Ribbon-Richtung (\hyperref[abk:L]{\texttt{'L'}}) der Wabenblöcke eindeutig zu markieren und parallel zur Prüf- bzw. Hauptlastrichtung einzubauen. Ein versehentlicher Quereinbau (\hyperref[abk:W]{\texttt{'W'}}) reduziert den Kernschubmodul von $550\,\text{MPa}$ auf $230\,\text{MPa}$ und halbiert die Schubtragfähigkeit.
\end{itemize}

\section{Einfluss der Kerndicke \& Parameterstudien} \label{sec:kerndicke}

Die Biegesteifigkeit (\hyperref[abk:EI]{$EI$}) eines Sandwich-Panels wird primär durch den Abstand der Deckhäute bestimmt (Steiner-Anteil):
\begin{equation}
    (EI)_{\text{equiv}} \approx E_{\text{skin}} \cdot b \cdot t_{\text{skin}} \cdot \frac{d^2}{2}
\end{equation}
wobei $d = t_{\text{core}} + t_{\text{skin}}$ den Hebelarm der Deckhautmitten darstellt.

\subsection*{Kerndicken im direkten Vergleich}
\begin{itemize}
    \item \textbf{10\,mm Kern}: Sehr kompakter Aufbau für enge Bauraumbereiche. Wegen des geringen Abstands treten jedoch sehr hohe Kernschubspannungen auf ($\tau_{\text{core}} > 2{,}4\,\text{MPa}$). Bei geringen Lagenzahlen dominiert Kernschubbruch als primärer Versagensmodus.
    \item \textbf{18\,mm Kern}: Der bewährte Standard mit optimaler Balance aus Bauraumbedarf, Beulsteifigkeit und Kernschubreserve.
    \item \textbf{20\,mm Kern}: Durch die Erhöhung von $18$ auf $20\,\text{mm}$ wächst der Hebelarm $d$ um über $11\,\%$. Dadurch kann die Deckhautdicke bei gleichem Biegegradienten signifikant reduziert werden.
\end{itemize}

  \chapter{Referenz: Materialkatalog \& \texorpdfstring{\abk{FSG}}{FSG}-Regeln} \label{ch:materials}

Dieses Kapitel dokumentiert alle in der zentralen Materialdatenbank (\hyperref[abk:CSV]{\texttt{materials.csv}}) hinterlegten Werkstoffe, deren Kennwerte und anerkannte Syntax-Aliase sowie die maßgeblichen Laminat-Regeln nach dem Formula-Student-Regelwerk.

\section{Material-Datenbank \& Kennwerte} \label{sec:materials_db}

Alle Materialkarten werden über die Datei \texttt{02\_Data/materials.csv} als Single Source of Truth verwaltet. Die Werte entsprechen den mechanischen Eigenschaften nach der Aushärtung im Heißpressverfahren ($\approx 6\,\text{bar}$):

\begin{table}[htbp]
\centering
\footnotesize
\setlength{\tabcolsep}{4.5pt}
\begin{tabular}{lllrrrrl}
\toprule
\textbf{Token} & \textbf{Bezeichnung} & \textbf{Typ} & \boldmath$E_1$ & \boldmath$X_t$ & \boldmath$t_{\text{ply}}$ & \boldmath$\rho$ & \textbf{Aliase} \\
 & & & \scriptsize [GPa] & \scriptsize [MPa] & \scriptsize [mm] & \scriptsize [$\text{g/cm}^3$] & \\
\midrule
\texttt{f}  & M40J Fabric     & \abk{CFK}-Gewebe   & $86{,}8$  & $864$  & $0{,}23$ & $1{,}49$ & \texttt{fm40j}, \texttt{m40jf}, \texttt{mf} \\
\texttt{m}  & M40J UD         & \abk{CFK}-\abk{UD}       & $175{,}6$ & $2273$ & $0{,}21$ & $1{,}52$ & \texttt{ud}, \texttt{m40jud} \\
\texttt{t7} & T700 Fabric     & \abk{CFK}-Gewebe   & $58{,}2$  & $1120$ & $0{,}31$ & $1{,}51$ & \texttt{t}, \texttt{t700}, \texttt{ft} \\
\texttt{t8} & T800S UD        & \abk{CFK}-\abk{UD}       & $155{,}7$ & $3344$ & $0{,}15$ & $1{,}54$ & \texttt{t800}, \texttt{t800s}, \texttt{u8} \\
\texttt{z}  & ZZ300T          & Zylon-Gewebe & $51{,}9$  & $988$  & $0{,}38$ & $0{,}78$ & \texttt{zz300t} \\
\texttt{a}  & SIGAPREG A W173 & Aramid-Gewebe& $26{,}0$  & $590$  & $0{,}23$ & $1{,}33$ & \texttt{aramid}, \texttt{kevlar} \\
\texttt{cp} & Coupling Film   & Klebefilm (\abk{CP})    & --        & --     & $0{,}15$ & $1{,}00$ & \texttt{cp} \\
\texttt{core}& Aluminium 5056 & Wabenkern    & --        & --     & $18{,}0$ & $0{,}075$& \texttt{core}, \texttt{corew} \\
\bottomrule
\end{tabular}
\caption{Übersicht der hinterlegten Werkstoffdaten in \texttt{materials.csv}.}
\label{tab:material_database}
\end{table}

\subsection*{Einsatzzwecke der Kernmaterialien}
\begin{itemize}
    \item \textbf{M40J \abk{UD} (\texttt{m})}: \abklink{HM}{Hochmodulfaser} für maximale Längssteifigkeit ($E_1 = 175{,}6\,\text{GPa}$). Bildet das Rückgrat für hohe Biegegradienten bei minimaler Lagenzahl.
    \item \textbf{T700 Fabric (\texttt{t7})}: Hochfestes Standardgewebe ($X_t = 1120\,\text{MPa}$). Durch die hohe Dehnfähigkeit ideal als robuste, schlagfeste Sichtlage gemäß \abk{FSG}-Reglement.
    \item \textbf{T800S \abk{UD} (\texttt{t8})}: \abklink{IM}{Intermediat-Modul-Faser} mit extremer Zugfestigkeit ($X_t = 3344\,\text{MPa}$). Eignet sich besonders für zugbeanspruchte Unterseiten zur Verhinderung vorzeitigen Faserbruchs.
    \item \textbf{M40J Fabric (\texttt{f})}: \abklink{HM}{Hochmodul}-Gewebe ($E_1 = 86{,}8\,\text{GPa}$). Verbindet hohe Steifigkeit mit Gewebestruktur; optimal als Anbindelage am Kern.
    \item \textbf{Zylon ZZ300T (\texttt{z})}: \abk{PBO}-Faser mit extremer Bruchzähigkeit und Energieaufnahme für Crash- und Splitterschutzzonen.
    \item \textbf{Klebefilm (\abk{CP})}: Unverzichtbare Epoxidharz-Kopplungsschicht ($\approx 150\,\text{g/m}^2$) zur formschlüssigen Wabenkernanbindung.
\end{itemize}

\section{Laminat-Regeln nach Formula Student (\texorpdfstring{\abk{FSG}}{FSG} T 3.5)} \label{sec:fsg_rules}

Bei der Auslegung und manuellen Definition von Laminaten müssen folgende Regelwerksgrenzen zwingend eingehalten werden:

\begin{itemize}
    \item \textbf{Sichtlagen-Regel (\abk{FSG} T\,3.5.7)}: \par
    Die äußerste Lage jeder Sandwich-Deckhaut muss zwingend einlagiges Gewebe sein (z.\,B. \texttt{45t7} oder \texttt{45f}). Reines \abk{UD}-Gelege an der Außenseite ist nach Regelwerk unzulässig.

    \item \textbf{Parallellagen-Begrenzung (\abk{FSG} T 3.5.3)}: \par
    Maximal $50\,\%$ der Einzellagen einer Deckhaut dürfen in die gleiche Faserrichtung orientiert sein. Beispielsweise sind bei einem 4-Lagen-Aufbau höchstens zwei $0^\circ$-Lagen erlaubt.

    \item \textbf{Asymmetrie-Regel der Deckhäute (\abk{FSG} T 3.5.4)}: \par
    Unterscheiden sich die Dicken der oberen und unteren Deckhaut ($t_{\text{top}} \ne t_{\text{lower}}$), darf das Dickenverhältnis nicht beliebig gewählt werden. Die dünnere Deckhaut ($t_{\text{min}}$) muss mindestens $40\,\%$ der Dicke der stärkeren Deckhaut ($t_{\text{max}}$) betragen:
    \begin{equation}
        t_{\text{min}} \ge \min(0{,}40 \cdot t_{\text{max}}, \, 1{,}0\,\text{mm})
    \end{equation}
    Erreicht die dünnere Deckhaut mindestens $1{,}0\,\text{mm}$ Dicke, ist jedes Dickenverhältnis regelkonform. Beim physischen 3-Punkt-Biegeversuch für den \abk{SES}-Nachweis muss bei asymmetrischen Aufbauten die schwächere Deckhaut auf Zug (Unterseite) getestet werden.

\end{itemize}


  \end{document}
